\section{Levantamiento y Análisis de Requerimientos}
\subsubsection{Requerimientos Funcionales}

% -----------------------------------------------------------

Los requerimientos funcionales definen las capacidades que el sistema debe proveer
para satisfacer los objetivos espec\'ificos (OE-01 a OE-13). Se organizan en los
mismos ocho m\'odulos definidos en el Cap\'itulo~1 y se presentan en la
Tabla~\ref{tab:rf}.

\setstretch{1.2}
\setlength{\tabcolsep}{4pt}
\begin{longtable}{>{\bfseries\centering\arraybackslash}p{1.5cm}
                >{\raggedright\arraybackslash}p{6.5cm}
                >{\raggedright\arraybackslash}p{6.5cm}
                >{\centering\arraybackslash}p{1.5cm}}
\caption{Requerimientos Funcionales del Sistema} \label{tab:rf} \\
\toprule
\thead{ID} & \thead{Requerimiento} & \thead{Criterio de aceptación} & \thead{OE ref.} \\
\midrule
\endfirsthead
\caption[]{TABLA \thetable{} (Continuación)} \\
\toprule
\thead{ID} & \thead{Requerimiento} & \thead{Criterio de aceptación} & \thead{OE ref.} \\
\midrule
\endhead
\bottomrule
\endfoot
\multicolumn{4}{l}{\textit{Módulo 1 — Infraestructura y Estructura Organizacional}} \\
\midrule
RF-01 & El sistema debe desplegar ambientes diferenciados de Producci\'on y Preview con pipelines de CI/CD autom\'aticos. &
  Deploy autom\'atico al hacer push a la rama \texttt{main}. Migraciones de Prisma ejecutadas sin errores en ambos ambientes. & OE-01 \\
\rowcolor{altrow}
RF-02 & El sistema debe gestionar los tres niveles organizacionales: Cadena, Zona y Sucursal, con integridad referencial. &
  CRUD completo en cada nivel. Intento de eliminar una Cadena con Zonas activas retorna error de integridad. & OE-02 \\
\midrule
\multicolumn{4}{l}{\textit{Módulo 2 — Autenticación y Control de Acceso}} \\
\midrule
RF-03 & El sistema debe autenticar usuarios colaboradores mediante email/contrase\~na y emitir JWT con el rol asignado. &
  Login exitoso genera token con claims de rol y nivel org. Token inv\'alido retorna HTTP 401. & OE-03 \\
\rowcolor{altrow}
RF-04 & Las pol\'iticas RLS deben garantizar que cada usuario acceda s\'olo a los datos de su nivel organizacional. &
  Gerente de Zona A no puede leer ni modificar registros de Zona B. Verificado v\'ia tests de integraci\'on. & OE-03 \\
RF-05 & Al primer login con contrase\~na temporal, el sistema debe forzar el cambio de contrase\~na antes de continuar. &
  El sistema redirige al formulario de cambio si el flag \texttt{first\_login = true}. Imposible omitir el paso. & OE-03 \\
\rowcolor{altrow}
RF-06 & Un usuario no puede asignar roles superiores a los propios al dar de alta nuevos colaboradores. &
  Gerente de Sucursal que intenta crear un Gerente de Zona recibe HTTP~403. & OE-04 \\
\midrule
\multicolumn{4}{l}{\textit{Módulo 3 — Catálogo}} \\
\midrule
RF-07 & El sistema debe permitir crear, editar, activar y desactivar categor\'ias y productos del men\'u por sucursal. &
  Producto con 3 variantes creado exitosamente. Producto inactivo no visible para comensales. & OE-05 \\
\rowcolor{altrow}
RF-08 & Los cambios al men\'u deben reflejarse para los comensales activos en menos de 2 segundos. &
  Test: Gerente desactiva producto; comensal conectado no lo ve en $< 2$~s sin recargar. & OE-05 \\
\midrule
\multicolumn{4}{l}{\textit{Módulo 4 — Mesas Físicas y Virtuales}} \\
\midrule
RF-09 & El sistema debe mostrar en tiempo real el estatus de todas las mesas f\'isicas de la sucursal. &
  Mesero ve el cambio de estatus de una mesa en $< 2$~s v\'ia Supabase Realtime sin recargar la p\'agina. & OE-06 \\
\rowcolor{altrow}
RF-10 & El sistema debe generar c\'odigos QR con token JWT de un solo uso (TTL 30~min) para la activaci\'on de mesas virtuales. &
  Segundo intento de escaneo del mismo QR retorna error de token ya utilizado o expirado. & OE-07 \\
RF-11 & Solo se pueden asociar mesas f\'isicas en estado LIBRE a una mesa virtual nueva (RB-M01). &
  Intento de inicializar con mesa en estado POR\_LIQUIDAR retorna error descriptivo. & OE-07 \\
\rowcolor{altrow}
RF-12 & El anfitri\'on debe crear una contrase\~na para activar la mesa virtual; esta se almacena hasheada (bcrypt). &
  Hash verificado en base de datos. Contrase\~na no recuperable en texto plano. & OE-07 \\
RF-13 & El comensal puede unirse a una mesa activa ingresando el ID de mesa y la contrase\~na provista por el anfitri\'on. &
  Credenciales incorrectas retornan HTTP~401. Credenciales correctas crean sesi\'on activa. & OE-07 \\
\midrule
\multicolumn{4}{l}{\textit{Módulo 5 — Órdenes y Cocina}} \\
\midrule
\rowcolor{altrow}
RF-14 & El comensal puede agregar productos a su canasta, ajustar cantidades y confirmar la orden con di\'alogo expl\'icito. &
  Canasta persiste durante la sesi\'on (Zustand). Confirmaci\'on sin di\'alogo: bloqueada. & OE-08 \\
RF-15 & El sistema debe distinguir entre \'ordenes individuales (cargo a cuenta personal) y \'ordenes compartidas (cargo a mesa grupal). &
  Orden individual aparece s\'olo en la cuenta del comensal. Orden compartida visible para todos los de la mesa. & OE-08 \\
\rowcolor{altrow}
RF-16 & El tablero del Chef debe mostrar \'ordenes nuevas en $< 2$~s con notificaci\'on sonora (Web Audio API). &
  Test: comensal confirma orden; aparece en tablero del Chef en $< 2$~s con sonido. & OE-09 \\
RF-17 & El Chef debe poder avanzar el estado de cada orden: PENDIENTE $\rightarrow$ EN\_PREPARACION $\rightarrow$ EN\_REPARTO. &
  Transici\'on inversa bloqueada. Cambios reflejados en vista del comensal en tiempo real. & OE-09 \\
\midrule
\multicolumn{4}{l}{\textit{Módulo 6 — Liquidación}} \\
\midrule
\rowcolor{altrow}
RF-18 & El sistema debe prevenir pagos concurrentes duplicados mediante \texttt{SELECT FOR UPDATE} en PostgreSQL. &
  Test de dos peticiones simult\'aneas: solo una se procesa; la segunda recibe error descriptivo. & OE-10 \\
RF-19 & La mesa virtual debe transicionar autom\'aticamente a LIQUIDADA cuando todos los comensales y la mesa compartida est\'en saldados. &
  Trigger en \texttt{transacciones\_pago}: al liquidar el \'ultimo pendiente, \texttt{mesa\_virtual.estado} = LIQUIDADA. & OE-10 \\
\rowcolor{altrow}
RF-20 & El sistema debe emitir un comprobante digital de liquidaci\'on al completarse todos los pagos de una mesa. &
  Comprobante generado con timestamp, \'items y montos. Verificado en ambiente de staging. & OE-11 \\
\midrule
\multicolumn{4}{l}{\textit{Módulo 7 — Analytics}} \\
\midrule
RF-21 & El dashboard debe filtrar m\'etricas autom\'aticamente seg\'un el nivel organizacional del usuario autenticado. &
  Admin ve toda la cadena; Gerente de Zona ve solo su zona; Gerente de Sucursal solo su sucursal. & OE-12 \\
\rowcolor{altrow}
RF-22 & El dashboard debe mostrar: n\'umero de mesas activas, \'ordenes por per\'iodo, ingresos, top productos y tiempos de atenci\'on. &
  Datos coherentes con las transacciones registradas en Supabase. & OE-12 \\
\midrule
\multicolumn{4}{l}{\textit{Módulo 8 — Seguridad}} \\
\midrule
RF-23 & Todas las comunicaciones deben estar cifradas con HTTPS/TLS~1.2+ (Vercel + Supabase). &
  Escaneo OWASP~ZAP sin vulnerabilidades cr\'iticas de cifrado. & OE-13 \\
\rowcolor{altrow}
RF-24 & Los endpoints de autenticaci\'on y acceso a mesa virtual deben tener rate limiting (m\'ax.\ 5 intentos fallidos antes de bloqueo temporal). &
  Test de fuerza bruta: sexto intento retorna HTTP~429 con cabecera \texttt{Retry-After}. & OE-13 \\
\end{longtable}
\setstretch{1.5}

% -----------------------------------------------------------
\subsubsection{Requerimientos No Funcionales}
% -----------------------------------------------------------

Los requerimientos no funcionales establecen las condiciones de calidad bajo las
cuales el sistema debe operar. La Tabla~\ref{tab:rnf} los organiza por categor\'ia.

\setstretch{1.2}
\setlength{\tabcolsep}{2pt}
\begin{longtable}{>{\bfseries\centering\arraybackslash}p{1.2cm}
                >{\raggedright\arraybackslash}p{2.3cm}
                >{\raggedright\arraybackslash}p{5.8cm}
                >{\raggedright\arraybackslash}p{5.8cm}}
\caption{Requerimientos No Funcionales del Sistema} \label{tab:rnf} \\
\toprule
\thead{ID} & \thead{Categoría} & \thead{Requerimiento} & \thead{Métrica / Criterio} \\
\midrule
\endfirsthead
\caption[]{TABLA \thetable{} (Continuación)} \\
\toprule
\thead{ID} & \thead{Categoría} & \thead{Requerimiento} & \thead{Métrica / Criterio} \\
\midrule
\endhead
\bottomrule
\endfoot
RNF-01 & Rendimiento & Tiempo de carga de la interfaz principal en conexi\'on m\'ovil 4G. &
  $\leq 3$~segundos medido con Lighthouse en modo m\'ovil. \\
\rowcolor{altrow}
RNF-02 & Rendimiento & Latencia de actualizaciones en tiempo real (\'ordenes, mesas, pagos). &
  $< 2$~segundos extremo a extremo v\'ia Supabase Realtime \cite{supabase_license}. \\
RNF-03 & Rendimiento & Soporte de accesos simult\'aneos en horas pico por sucursal. &
  Sistema estable con hasta 200 conexiones concurrentes. Verificado con k6. \\
\rowcolor{altrow}
RNF-04 & Disponibilidad & Disponibilidad durante el horario operativo del restaurante. &
  $\geq 99.5$\% de uptime mensual (SLA de Vercel Pro + Supabase Pro) \cite{vercel_pricing, supabase_pricing}. \\
RNF-05 & Seguridad & Transmisi\'on cifrada en todas las rutas de la plataforma. &
  HTTPS/TLS~1.2+ en Vercel y Supabase. SSL gestionado autom\'aticamente. \\
\rowcolor{altrow}
RNF-06 & Seguridad & Protecci\'on de datos sensibles en base de datos. &
  RLS habilitado en todas las tablas sensibles. Sin datos de tarjeta almacenados. \\
RNF-07 & Seguridad & Integridad de datos transaccionales bajo concurrencia. &
  Transacciones ACID with \texttt{SELECT FOR UPDATE}. Sin cobros duplicados verificados
  en pruebas de carga. \\
\rowcolor{altrow}
RNF-08 & Usabilidad & Interfaz Mobile-First adaptable a pantallas $\geq 375$~px. &
  Responsive verificado en Chrome DevTools para iPhone~SE, Galaxy~S20 y iPad. \\
RNF-09 & Usabilidad & Accesibilidad WCAG~2.1~AA: contraste $\geq 4.5$:1 para texto. &
  Auditado con axe DevTools. Sin errores cr\'iticos de accesibilidad. \\
\rowcolor{altrow}
RNF-10 & Usabilidad & Curva de aprendizaje del personal operativo $\leq 4$ horas. &
  Onboarding guiado interactivo en la app. Manual de usuario por rol. \\
RNF-11 & Compatibilidad & Compatibilidad con navegadores modernos sin instalaci\'on nativa. &
  Funcional en Chrome~$\geq$90, Safari~$\geq$14 y Firefox~$\geq$88. \\
\rowcolor{altrow}
RNF-12 & Escalabilidad & Arquitectura con escalado horizontal autom\'atico. &
  Funciones serverless en Vercel escalan autom\'aticamente.
  PgBouncer (Supabase) resuelve pool de conexiones \cite{mdn_pwa, ibm_bi_intro}. \\
RNF-13 & Mantenibilidad & C\'odigo base con cobertura de pruebas unitarias $\geq 70$\%. &
  Reportado por Jest/Vitest en pipeline de CI/CD al hacer merge a \texttt{main}. \\
\rowcolor{altrow}
RNF-14 & Mantenibilidad & Tipado est\'atico en todo el stack. &
  TypeScript strict mode habilitado. Sin \texttt{any} impl\'icito en c\'odigo de producci\'on. \\
RNF-15 & Legal & Cumplimiento de LFPDPPP antes del primer usuario real. &
  Aviso de Privacidad publicado. Consentimiento expl\'icito implementado.
  Mecanismo ARCO disponible. \\
\end{longtable}
\setstretch{1.5}



% ============================================================
\newpage
% Casos de Uso (movidos a Cap. 4)
\subsection{Diagramas de Casos de Uso UML}
% -----------------------------------------------------------

Los diagramas de casos de uso modelan las interacciones entre los actores del
sistema y sus funcionalidades, siguiendo la notaci\'on UML~2.5.1 \cite{omg_uml_2017}.
Se definen quince casos de uso principales (CU-01 a CU-15), consistentes con los
objetivos espec\'ificos OE-06 a OE-12 y los requerimientos funcionales RF-01 a
RF-22 descritos en la Secci\'on~4.2. Cada diagrama presenta los actores
involucrados, las relaciones de inclusi\'on (\textit{include}), extensi\'on
(\textit{extend}) y herencia, y la frontera del sistema.

\begingroup
\setstretch{1.0}
\renewcommand{\arraystretch}{0.95}
\setlength{\tabcolsep}{2pt}
\subsubsection{CU-01: Comensal Escanea QR e Ingresa a Mesa Virtual}

La Figura~\ref{fig:cu01} presenta el caso de uso del flujo de acceso a la mesa virtual mediante
c\'odigo QR. El actor principal es el \textbf{Comensal} y el actor secundario el
\textbf{Mesero}, quien inicializa la mesa. El sistema valida el token JWT del QR
(RF-10) y ejecuta el flujo de activaci\'on descrito en la Figura~\ref{fig:seq_activacion}.

\begin{figure}[H]
\centering
\includegraphics[width=0.68\textwidth]{casos_uso/rendered/CU_01.pdf}
\vspace{2pt}
\scriptsize
\begin{tabular}{>{\raggedright\arraybackslash}p{2.6cm} >{\raggedright\arraybackslash}p{11.8cm}}
\textbf{Actores} & Mesero (primario), Comensal / Anfitr\'ion (primario), Sistema \\
\textbf{Precondiciones} & Mesa f\'isica en estado \texttt{DISPONIBLE}.
  El mesero ha iniciado sesi\'on con rol MESERO o superior. \\
\textbf{Flujo principal} &
  \textbf{1.} Mesero selecciona mesas f\'isicas y genera mesa virtual.\newline
  \textbf{2.} Sistema genera QR con JWT firmado (TTL 30 min) --- \textit{include} \textbf{CU-01a}: Generar token JWT de un solo uso.\newline
  \textbf{3.} Mesero entrega QR f\'isico o digital al anfitr\'ion del grupo.\newline
  \textbf{4.} Anfitr\'ion escanea el QR desde su navegador m\'ovil.\newline
  \textbf{5.} Sistema valida firma, TTL y estado del token.\newline
  \textbf{6.} Anfitr\'ion ingresa nombre y crea contrase\~na de mesa.\newline
  \textbf{7.} Sistema activa la mesa virtual (\texttt{PENDIENTE\_ACTIVACION} $\to$ \texttt{OCUPADA}).\newline
  \textbf{8.} Anfitr\'ion recibe \texttt{joinCode} para compartir con sus compa\~neros.\newline
  \textbf{9.} Comensales adicionales ingresan con ID de mesa + contrase\~na. \\
\textbf{Flujos alternativos} &
  \textit{5a.} Token expirado o ya utilizado $\to$ Sistema retorna error \texttt{HTTP 401}.\newline
  \textit{5b.} Mesa ya activa $\to$ Sistema retorna error \texttt{HTTP 409}.\newline
  \textit{6a.} Contrase\~na de menos de 4 caracteres $\to$ validaci\'on en frontend. \\
\textbf{Postcondiciones} & Mesa virtual activa. Sesi\'on de Anfitr\'ion creada con \texttt{isHost=true}.
  Mesas f\'isicas asociadas en estado \texttt{OCUPADA}. \\
\textbf{RF relacionados} & RF-09, RF-10, RF-11, RF-12, RF-13 \\
\end{tabular}
\caption{CU-01: Comensal escanea QR e ingresa a mesa virtual. Relaci\'on de inclusi\'on con CU-01a (generaci\'on de token JWT de un solo uso).}
\label{fig:cu01}
\end{figure}

\newpage

\subsubsection{CU-02: Comensal Consulta Menú, Crea Pedido y Confirma Orden}

La Figura~\ref{fig:cu02} modela el flujo de creaci\'on de orden desde la perspectiva del
comensal. Incluye la consulta del men\'u digital, la gesti\'on de la canasta
y la confirmaci\'on expl\'icita de la orden.

\begin{figure}[H]
\centering
\includegraphics[width=0.68\textwidth]{casos_uso/rendered/CU_02.pdf}
\vspace{2pt}
\scriptsize
\begin{tabular}{>{\raggedright\arraybackslash}p{2.6cm} >{\raggedright\arraybackslash}p{11.8cm}}
\textbf{Actores} & Comensal (primario), Sistema, Chef (actor receptor) \\
\textbf{Precondiciones} & Comensal con sesi\'on activa en una mesa virtual (\texttt{isActive=true}). \\
\textbf{Flujo principal} &
  \textbf{1.} Comensal navega el men\'u digital por categor\'ias.\newline
  \textbf{2.} Comensal agrega productos a su canasta (Optimistic UI).\newline
  \textbf{3.} Comensal selecciona tipo de orden: \textbf{Individual} (cargo a su cuenta) o \textbf{Compartida} (cargo a la cuenta de la mesa) --- \textit{extend} \textbf{CU-02a}: Agregar variante de producto.\newline
  \textbf{4.} Comensal confirma la orden mediante di\'alogo expl\'icito.\newline
  \textbf{5.} Sistema crea \texttt{Order} y \texttt{OrderItem}s at\'omicamente y actualiza la \texttt{Cuenta}.\newline
  \textbf{6.} Sistema emite evento \texttt{nueva\_orden} v\'ia Supabase Realtime.\newline
  \textbf{7.} Chef recibe la orden en el tablero KDS (CU-03). \\
\textbf{Flujos alternativos} &
  \textit{4a.} Mesa en estado \texttt{CERRANDO} $\to$ Sistema bloquea nuevas \'ordenes (\texttt{HTTP 409}).\newline
  \textit{2a.} Producto \texttt{isSoldOut=true} $\to$ No mostrado en el men\'u (RF-08). \\
\textbf{Postcondiciones} & Orden en estado \texttt{PENDING}. Cuenta del comensal actualizada. Orden visible en KDS del Chef. \\
\textbf{RF relacionados} & RF-07, RF-08, RF-14, RF-15 \\
\end{tabular}
\caption{CU-02: Comensal consulta men\'u, crea pedido y confirma orden. Relaci\'on de extensi\'on con CU-02a (selecci\'on de variante de producto).}
\label{fig:cu02}
\end{figure}

\subsubsection{CU-03: Cocina Recibe Orden, Actualiza Estado y Notifica Avance}

La Figura~\ref{fig:cu03} describe el flujo del Chef al gestionar \'ordenes en el KDS Kanban.
El flujo de actualizaci\'on de estado se sincroniza en tiempo real con la vista
del comensal mediante Supabase Realtime.

\begin{figure}[H]
\centering
\includegraphics[width=0.68\textwidth]{casos_uso/rendered/CU_03.pdf}
\vspace{2pt}
\scriptsize
\begin{tabular}{>{\raggedright\arraybackslash}p{2.6cm} >{\raggedright\arraybackslash}p{11.8cm}}
\textbf{Actores} & Chef (primario), Sistema, Comensal (actor receptor) \\
\textbf{Precondiciones} & Orden en estado \texttt{PENDING} en el tablero KDS. Chef autenticado con rol CHEF. \\
\textbf{Flujo principal} &
  \textbf{1.} Sistema notifica al Chef con alerta sonora (Web Audio API) al llegar una orden nueva.\newline
  \textbf{2.} Chef visualiza la orden en la columna ``Nuevas'' del tablero Kanban.\newline
  \textbf{3.} Chef filtra \'ordenes por estaci\'on (\texttt{COCINA} o \texttt{BARRA}) --- \textit{include} \textbf{CU-03a}: Filtrar \'ordenes por estaci\'on.\newline
  \textbf{4.} Chef avanza la orden: \texttt{PENDING} $\to$ \texttt{IN\_PROGRESS}.\newline
  \textbf{5.} Chef avanza la orden: \texttt{IN\_PROGRESS} $\to$ \texttt{READY}.\newline
  \textbf{6.} Sistema emite evento \texttt{orden\_actualizada} v\'ia Supabase Realtime.\newline
  \textbf{7.} Vista del comensal actualiza el estado de la orden en tiempo real.\newline
  \textbf{8.} Mesero entrega el platillo y avanza a \texttt{DELIVERED}. \\
\textbf{Flujos alternativos} &
  \textit{4a.} SLA $> 60$~s sin confirmaci\'on $\to$ Alerta al gerente (RO-05).\newline
  \textit{5a.} Chef marca orden err\'oneamente $\to$ Bot\'on \textit{Deshacer} disponible hasta siguiente transici\'on. \\
\textbf{Postcondiciones} & Orden en estado \texttt{DELIVERED}. Timestamp \texttt{deliveredAt} registrado. Vista del comensal actualizada. \\
\textbf{RF relacionados} & RF-16, RF-17 \\
\end{tabular}
\caption{CU-03: Cocina recibe orden, actualiza estado y notifica avance. Inclusi\'on con CU-03a (filtrado por estaci\'on de preparaci\'on).}
\label{fig:cu03}
\end{figure}

\subsubsection{CU-04: Comensal Liquida Cuenta (Individual o Mesa)}

La Figura~\ref{fig:cu04} detalla el proceso de cierre y pago de cuentas. Este
flujo es cr\'itico para la integridad financiera del sistema, integrando el
bloqueo de concurrencia (RF-18) y la transici\'on de estado de la mesa (RF-19).

\begin{figure}[H]
\centering
\includegraphics[width=0.68\textwidth]{casos_uso/rendered/CU_04.pdf}
\vspace{2pt}
\scriptsize
\begin{tabular}{>{\raggedright\arraybackslash}p{2.6cm} >{\raggedright\arraybackslash}p{11.8cm}}
\textbf{Actores} & Comensal (primario), Sistema, Pasarela de Pagos (externo) \\
\textbf{Precondiciones} & Mesa virtual en estado \texttt{OCUPADA} o \texttt{POR\_LIQUIDAR}.
  Comensal con ítems por pagar. \\
\textbf{Flujo principal} &
  \textbf{1.} Comensal selecciona ``Pagar Cuenta'' en su interfaz.\newline
  \textbf{2.} Comensal elige entre liquidar su cuenta individual o la compartida de la mesa.\newline
  \textbf{3.} Sistema bloquea los registros de la cuenta (\texttt{SELECT FOR UPDATE}) para evitar pagos duplicados --- \textit{include} \textbf{CU-04a}: Bloquear concurrencia de pago.\newline
  \textbf{4.} Comensal ingresa datos de pago y confirma transacci\'on.\newline
  \textbf{5.} Sistema valida con pasarela externa y marca cuenta como \texttt{PAID}.\newline
  \textbf{6.} Sistema emite evento \texttt{pago\_confirmado} v\'ia Supabase Realtime.\newline
  \textbf{7.} Si la cuenta de la mesa llega a saldo cero, el sistema marca la mesa como \texttt{POR\_LIBERAR}.\newline
  \textbf{8.} Mesero confirma liberaci\'on f\'isica de la mesa (RF-19). \\
\textbf{Flujos alternativos} &
  \textit{3a.} Registro bloqueado por otro pago en curso $\to$ Sistema solicita esperar.\newline
  \textit{5a.} Transacci\'on rechazada $\to$ Sistema notifica motivo y permite reintentar.\newline
  \textit{7a.} Saldo pendiente remanente $\to$ Mesa permanece en \texttt{OCUPADA}. \\
\textbf{Postcondiciones} & Cuenta liquidada. Transacci\'on registrada en el historial. Mesa liberada (si aplica). \\
\textbf{RF relacionados} & RF-18, RF-19, RF-20 \\
\end{tabular}
\caption{CU-04: Comensal liquida cuenta (individual o mesa). Incluye bloqueo de registros para control de concurrencia.}
\label{fig:cu04}
\end{figure}

\subsubsection{CU-05: Gerente Consulta Dashboard de Analítica}

La Figura~\ref{fig:cu05} ilustra el flujo de consulta de inteligencia de negocio.
El sistema garantiza que la informaci\'on sea filtrada seg\'un el nivel de
autoridad del usuario (RF-21).

\begin{figure}[H]
\centering
\includegraphics[width=0.68\textwidth]{casos_uso/rendered/CU_05.pdf}
\vspace{2pt}
\scriptsize
\begin{tabular}{>{\raggedright\arraybackslash}p{2.6cm} >{\raggedright\arraybackslash}p{11.8cm}}
\textbf{Actores} & Gerente de Sucursal / Zona / Cadena (primario), Sistema \\
\textbf{Precondiciones} & Gerente con sesi\'on activa y permisos de lectura en el m\'odulo de anal\'itica. \\
\textbf{Flujo principal} &
  \textbf{1.} Gerente accede al m\'odulo de "Analytics" en el dashboard administrativo.\newline
  \textbf{2.} Sistema consulta las tablas Gold de la base de datos (\texttt{analytic\_sales\_daily}, \texttt{analytic\_item\_velocity}, etc.).\newline
  \textbf{3.} Sistema renderiza m\'etricas clave: Ingresos Totales, Ticket Promedio, Productos Top.\newline
  \textbf{4.} Gerente filtra datos por rango de fechas y categor\'ia de producto.\newline
  \textbf{5.} Gerente visualiza el pron\'ostico de demanda para los pr\'oximos 7 d\'ias (generado por Spark).\newline
  \textbf{6.} Gerente identifica cuellos de botella mediante m\'etricas de \textit{Process Mining}. \\
\textbf{Postcondiciones} & Visualización de métricas de negocio actualizada. Soporte a la toma de decisiones basada en datos. \\
\textbf{RF relacionados} & RF-21, RF-22 \\
\end{tabular}
\caption{CU-05: Gerente consulta dashboard de anal\'itica. Visualizaci\'on de resultados precalculados por Spark.}
\label{fig:cu05}
\end{figure}

\newpage

\subsubsection{CU-06: Gestión de Estructura Organizacional y Personal}

La Figura~\ref{fig:cu06} describe el proceso de configuraci\'on de la jerarqu\'ia
empresarial y el control de acceso del personal, fundamental para la
seguridad multitenant (RF-01 a RF-06).

\begin{figure}[H]
\centering
\includegraphics[width=0.68\textwidth]{casos_uso/rendered/CU_06.pdf}
\vspace{2pt}
\scriptsize
\begin{tabular}{>{\raggedright\arraybackslash}p{2.6cm} >{\raggedright\arraybackslash}p{11.8cm}}
\textbf{Actores} & Administrador de Cadena (primario), Gerente (primario), Sistema \\
\textbf{Precondiciones} & Administrador autenticado con rol ADMIN o superior. \\
\textbf{Flujo principal} &
  \textbf{1.} Administrador accede a la configuraci\'on de "Organizaci\'on".\newline
  \textbf{2.} Administrador crea o edita Zonas y Sucursales vinculadas a la Cadena.\newline
  \textbf{3.} Administrador gestiona el personal (Staff), asignando roles y PINes de acceso.\newline
  \textbf{4.} Sistema valida unicidad de PINes dentro de la sucursal.\newline
  \textbf{5.} Sistema aplica Row Level Security (RLS) basado en el \texttt{orgId} del nuevo colaborador. \\
\textbf{Postcondiciones} & Estructura organizacional y personal actualizados en la base de datos. \\
\textbf{RF relacionados} & RF-01, RF-02, RF-03 \\
\end{tabular}
\caption{CU-06: Gestión de estructura organizacional y personal. Control de acceso jerárquico.}
\label{fig:cu06}
\end{figure}

\newpage

\subsubsection{CU-07: Administración de Catálogo y Overrides de Menú}

La Figura~\ref{fig:cu07} ilustra la gesti\'on flexible de productos, permitiendo
ajustes locales sin perder la centralizaci\'on de la cadena (RF-07, RF-14).

\begin{figure}[H]
\centering
\includegraphics[width=0.68\textwidth]{casos_uso/rendered/CU_07.pdf}
\vspace{2pt}
\scriptsize
\begin{tabular}{>{\raggedright\arraybackslash}p{2.6cm} >{\raggedright\arraybackslash}p{11.8cm}}
\textbf{Actores} & Gerente (primario), Sistema \\
\textbf{Precondiciones} & Gerente autenticado con permisos de gestión de catálogo. \\
\textbf{Flujo principal} &
  \textbf{1.} Gerente accede al "Gestor de Menú".\newline
  \textbf{2.} Sistema muestra la plantilla maestra (\textit{MenuTemplate}) heredada de la cadena.\newline
  \textbf{3.} Gerente aplica un \textit{Override} (sobreescritura) de precio o disponibilidad para un producto.\newline
  \textbf{4.} Sistema guarda el cambio en \texttt{RestaurantMenuOverride}.\newline
  \textbf{5.} Sistema invalida la cache de men\'u para esa sucursal. \\
\textbf{Postcondiciones} & Menú de la sucursal actualizado sin afectar a otros locales de la cadena. \\
\textbf{RF relacionados} & RF-04, RF-05, RF-06 \\
\end{tabular}
\caption{CU-07: Administración de catálogo y overrides de menú. Herencia y especialización por sucursal.}
\label{fig:cu07}
\end{figure}

\newpage

\subsubsection{CU-08: Diseño de Planta y Mapeo Visual de Mesas}

La Figura~\ref{fig:cu08} detalla la configuraci\'on de la infraestructura f\'isica
del restaurante, habilitando la interfaz visual del mesero (RF-09).

\begin{figure}[H]
\centering
\includegraphics[width=0.68\textwidth]{casos_uso/rendered/CU_08.pdf}
\vspace{2pt}
\scriptsize
\begin{tabular}{>{\raggedright\arraybackslash}p{2.6cm} >{\raggedright\arraybackslash}p{11.8cm}}
\textbf{Actores} & Gerente de Sucursal (primario), Sistema \\
\textbf{Precondiciones} & Gerente autenticado en la sucursal destino. \\
\textbf{Flujo principal} &
  \textbf{1.} Gerente accede al "Editor de Piso" (\textit{Floor Map Designer}).\newline
  \textbf{2.} Gerente arrastra y posiciona mesas en el lienzo interactivo (SVG/Canvas).\newline
  \textbf{3.} Gerente define el n\'umero y capacidad (pax) de cada mesa.\newline
  \textbf{4.} Sistema persiste las coordenadas $(x, y)$ y el \texttt{parentTableId} para uniones.\newline
  \textbf{5.} Sistema genera los c\'odigos QR est\'aticos para impresi\'on. \\
\textbf{Postcondiciones} & Plano visual de mesas actualizado y disponible para el equipo de meseros. \\
\textbf{RF relacionados} & RF-03, RF-12 \\
\end{tabular}
\caption{CU-08: Diseño de planta y mapeo visual de mesas. Gestión espacial de la operación.}
\label{fig:cu08}
\end{figure}

\newpage

\subsubsection{CU-09: Auditoría de Acciones y Trazabilidad de Cambios}

La Figura~\ref{fig:cu09} detalla el sistema de logs de auditor\'ia inmutables,
garantizando la transparencia en las operaciones administrativas y de cat\'alogo.

\begin{figure}[H]
\centering
\includegraphics[width=0.68\textwidth]{casos_uso/rendered/CU_09.pdf}
\vspace{2pt}
\scriptsize
\begin{tabular}{>{\raggedright\arraybackslash}p{2.6cm} >{\raggedright\arraybackslash}p{11.8cm}}
\textbf{Actores} & Administrador de Cadena (primario), Sistema \\
\textbf{Precondiciones} & Administrador con acceso al m\'odulo de seguridad. \\
\textbf{Flujo principal} &
  \textbf{1.} Administrador selecciona el m\'odulo de "Logs de Auditoría".\newline
  \textbf{2.} Sistema recupera registros de la tabla \texttt{AuditLog}.\newline
  \textbf{3.} Administrador filtra por usuario, tipo de acci\'on (ej. Reembolso) o fecha.\newline
  \textbf{4.} Sistema muestra el \textit{payload} JSON con los cambios realizados (antes/despu\'es). \\
\textbf{Postcondiciones} & Trazabilidad completa de acciones sensibles para prevención de fraudes. \\
\textbf{RF relacionados} & RF-01, RF-21 \\
\end{tabular}
\caption{CU-09: Auditoría de acciones y trazabilidad de cambios. Transparencia operativa.}
\label{fig:cu09}
\end{figure}

\newpage

\subsubsection{CU-10: Configuración de Modificadores y Opciones de Platillo}

La Figura~\ref{fig:cu10} describe la personalizaci\'on avanzada de productos,
necesaria para flujos gastron\'omicos complejos.

\begin{figure}[H]
\centering
\includegraphics[width=0.68\textwidth]{casos_uso/rendered/CU_10.pdf}
\vspace{2pt}
\scriptsize
\begin{tabular}{>{\raggedright\arraybackslash}p{2.6cm} >{\raggedright\arraybackslash}p{11.8cm}}
\textbf{Actores} & Gerente (primario), Sistema \\
\textbf{Precondiciones} & Producto existente en el catálogo. \\
\textbf{Flujo principal} &
  \textbf{1.} Gerente selecciona un producto en el dashboard administrativo.\newline
  \textbf{2.} Gerente define grupos de modificadores (ej. "Término de la carne").\newline
  \textbf{3.} Gerente asigna opciones a cada grupo, definiendo si tienen costo adicional.\newline
  \textbf{4.} Sistema asocia los modificadores al producto mediante \texttt{ProductModifierGroup}.\newline
  \textbf{5.} Comensal visualiza y selecciona las opciones al agregar el producto a su canasta. \\
\textbf{Postcondiciones} & Producto configurado con opciones personalizables. \\
\textbf{RF relacionados} & RF-14, RF-15 \\
\end{tabular}
\caption{CU-10: Configuración de modificadores y opciones de platillo. Flexibilidad en la personalización de órdenes.}
\label{fig:cu10}
\end{figure}

\newpage

\subsubsection{CU-11: Aplicación de Descuentos, Cortesías y Reembolsos}

La Figura~\ref{fig:cu11} modela el flujo de ajustes financieros y atenci\'on al
cliente por parte de la gerencia.

\begin{figure}[H]
\centering
\includegraphics[width=0.68\textwidth]{casos_uso/rendered/CU_11.pdf}
\vspace{2pt}
\scriptsize
\begin{tabular}{>{\raggedright\arraybackslash}p{2.6cm} >{\raggedright\arraybackslash}p{11.8cm}}
\textbf{Actores} & Gerente (primario), Sistema, Pasarela de Pagos (externo) \\
\textbf{Precondiciones} & Pago confirmado o cuenta abierta con ítems registrados. \\
\textbf{Flujo principal} &
  \textbf{1.} Gerente selecciona una cuenta o pago específico.\newline
  \textbf{2.} Gerente aplica un descuento (porcentaje o monto) o marca un ítem como cortesía.\newline
  \textbf{3.} Sistema recalcula el saldo total de la mesa en tiempo real.\newline
  \textbf{4.} En caso de reembolso, el sistema solicita la reversión a la pasarela externa.\newline
  \textbf{5.} Sistema registra la acci\'on en los logs de auditor\'ia con el motivo especificado. \\
\textbf{Postcondiciones} & Ajuste financiero aplicado y registrado para control interno. \\
\textbf{RF relacionados} & RF-18, RF-21 \\
\end{tabular}
\caption{CU-11: Aplicación de descuentos, cortesías y reembolsos. Gestión de excepciones financieras.}
\label{fig:cu11}
\end{figure}

\newpage

\subsubsection{CU-12: Gestión de Impuestos y Parámetros Fiscales}

La Figura~\ref{fig:cu12} describe la configuraci\'on de reg\'imenes fiscales por
cadena y sucursal.

\begin{figure}[H]
\centering
\includegraphics[width=0.68\textwidth]{casos_uso/rendered/CU_12.pdf}
\vspace{2pt}
\scriptsize
\begin{tabular}{>{\raggedright\arraybackslash}p{2.6cm} >{\raggedright\arraybackslash}p{11.8cm}}
\textbf{Actores} & Administrador de Cadena (primario), Sistema \\
\textbf{Precondiciones} & Configuración organizacional inicial completada. \\
\textbf{Flujo principal} &
  \textbf{1.} Administrador accede a "Parámetros Fiscales" en la configuración de la cadena.\newline
  \textbf{2.} Administrador define la tasa de IVA y otros impuestos locales por zona.\newline
  \textbf{3.} Administrador activa o desactiva la solicitud obligatoria de datos de facturación.\newline
  \textbf{4.} Sistema aplica estos parámetros en el cálculo automático de cada cuenta.\newline
  \textbf{5.} Sistema genera el desglose de impuestos en el comprobante digital. \\
\textbf{Postcondiciones} & Parámetros fiscales globales aplicados a todas las transacciones de la cadena. \\
\textbf{RF relacionados} & RF-18, RF-20 \\
\end{tabular}
\caption{CU-12: Gestión de impuestos y parámetros fiscales. Garantiza el cumplimiento normativo en las transacciones.}
\label{fig:cu12}
\end{figure}

\newpage

\subsubsection{CU-13: Ejecución de Análisis y Orquestación de Pipeline Medallion}

La Figura~\ref{fig:cu13} describe la interacci\'on entre el backend Next.js y
el servidor de anal\'itica Spark, implementando un patr\'on de orquestaci\'on as\'incrona.

\begin{figure}[H]
\centering
\includegraphics[width=0.68\textwidth]{casos_uso/rendered/CU_13.pdf}
\vspace{2pt}
\scriptsize
\begin{tabular}{>{\raggedright\arraybackslash}p{2.6cm} >{\raggedright\arraybackslash}p{11.8cm}}
\textbf{Actores} & Gerente (primario), Sistema (Next.js), Analytics Orchestrator (receptor) \\
\textbf{Precondiciones} & Al menos 24 horas de datos operacionales acumulados. \\
\textbf{Flujo principal} &
  \textbf{1.} Gerente solicita "Actualizar Análisis" desde el Dashboard.\newline
  \textbf{2.} Sistema valida permisos e invoca al \textit{Orchestrator} (FastAPI) en el VPS.\newline
  \textbf{3.} Orquestador registra el inicio en \texttt{AnalyticsJobRun} y dispara el job de Spark.\newline
  \textbf{4.} Dashboard inicia el \textit{polling} de estado mediante el \texttt{jobId}.\newline
  \textbf{5.} Spark procesa capas Bronze $\to$ Silver $\to$ Gold.\newline
  \textbf{6.} Orquestador marca el job como \texttt{COMPLETED}.\newline
  \textbf{7.} Dashboard detecta cambio de estado y refresca las gr\'aficas anal\'iticas. \\
\textbf{Postcondiciones} & Tablas analíticas Gold actualizadas en PostgreSQL. \\
\textbf{RF relacionados} & RF-21, RF-22 \\
\end{tabular}
\caption{CU-13: Ejecución de análisis y orquestación de pipeline Medallion. Separación de cargas OLTP y OLAP.}
\label{fig:cu13}
\end{figure}

\newpage

\subsubsection{CU-14: Generación de Pronósticos de Demanda (Demand Forecasting)}

La Figura~\ref{fig:cu14} modela la l\'ogica de inteligencia de negocio para la
predicci\'on de demanda basada en estacionalidad hist\'orica.

\begin{figure}[H]
\centering
\includegraphics[width=0.68\textwidth]{casos_uso/rendered/CU_14.pdf}
\vspace{2pt}
\scriptsize
\begin{tabular}{>{\raggedright\arraybackslash}p{2.6cm} >{\raggedright\arraybackslash}p{11.8cm}}
\textbf{Actores} & Sistema (Analytics Engine), Gerente (receptor) \\
\textbf{Precondiciones} & Datos hist\'oricos suficientes en capa Silver.\newline
\textbf{Flujo principal} &
  \textbf{1.} Spark extrae el hist\'orico de ventas de las \'ultimas 4 semanas.\newline
  \textbf{2.} Sistema calcula la media m\'ovil ponderada por d\'ia de la semana.\newline
  \textbf{3.} Sistema genera el \textit{Score de Confianza} basado en la variabilidad hist\'orica.\newline
  \textbf{4.} Sistema proyecta la demanda para los pr\'oximos 7 d\'ias.\newline
  \textbf{5.} Resultados se persisten en \texttt{analytic\_forecast} para su visualizaci\'on. \\
\textbf{Postcondiciones} & Proyecciones de venta disponibles en el dashboard administrativo. \\
\textbf{RF relacionados} & RF-21 \\
\end{tabular}
\caption{CU-14: Generación de pronósticos de demanda. Proporciona una visión predictiva para la toma de decisiones.}
\label{fig:cu14}
\end{figure}

\subsubsection{CU-15: Análisis de Tiempos de Servicio y Eficiencia (Process Mining)}

La Figura~\ref{fig:cu15} detalla la monitorizaci\'on de la eficiencia operativa
en cocina y sal\'on.

\begin{figure}[H]
\centering
\includegraphics[width=0.68\textwidth]{casos_uso/rendered/CU_15.pdf}
\vspace{2pt}
\scriptsize
\begin{tabular}{>{\raggedright\arraybackslash}p{2.6cm} >{\raggedright\arraybackslash}p{11.8cm}}
\textbf{Actores} & Sistema (Analytics Engine), Gerente (receptor), Chef (receptor) \\
\textbf{Precondiciones} & Eventos de estado de orden registrados en la base de datos. \\
\textbf{Flujo principal} &
  \textbf{1.} Spark analiza la cronolog\'ia de \texttt{OrderItemStatusEvent}.\newline
  \textbf{2.} Sistema calcula el tiempo promedio de ciclo (\texttt{NEW} $\to$ \texttt{SERVED}).\newline
  \textbf{3.} Sistema identifica cuellos de botella mediante el \texttt{bottleneckScore} por estaci\'on.\newline
  \textbf{4.} Sistema detecta desviaciones respecto al SLA configurado.\newline
  \textbf{5.} M\'etricas de eficiencia se despliegan en el m\'odulo de Operaciones. \\
\textbf{Postcondiciones} & Identificación de áreas de mejora operativa y saturación de cocina. \\
\textbf{RF relacionados} & RF-16, RF-17, RF-22 \\
\end{tabular}
\caption{CU-15: Análisis de tiempos de servicio y eficiencia. Permite el monitoreo continuo de la calidad del servicio.}
\label{fig:cu15}
\end{figure}


\endgroup
%  CAPÍTULO 5 — DISEÑO DEL SISTEMA
% ============================================================
\newpage
